% Unit 070 English translation of chapter5.tex lines 955--1063.
% Source: Methods of Algebra, Volume 2, commit
% 9a5803ff77dd3257484cb177f851a73770a59dd3 (CC BY 4.0).
% stable-unit-id: o014.aljabr2.chapter5.multiplicative-structures
% Coverage: Complete Section 5.7, before the exercises.

% segment-id: o014.li2.chapter5.multiplicative-structures.h001
\section{A glimpse of multiplicative structures}\label{sec:multiplicative-SS}
\hypertarget{o014.aljabr2.chapter5.multiplicative-structures}{ }

% segment-id: o014.li2.chapter5.multiplicative-structures.p001
For simplicity and concreteness, in this section we take a commutative ring
$\Bbbk$ and $\mathcal{A}=\Bbbk\dcate{Mod}$; this suffices for the classical
applications. We refer to $\Bbbk$-modules simply as modules and to
$\Bbbk$-algebras simply as algebras, and write
$\otimes:=\otimes_{\Bbbk}$.

% segment-id: o014.li2.chapter5.multiplicative-structures.p002
Let $I$ be a commutative monoid, with its binary operation written additively.
The definitions from \cite[\S 7.4]{Li1} may be summarized as follows:
\begin{itemize}
	\item an $I$-graded module is a module with a direct-sum decomposition
	$M=\bigoplus_{i\in I}M^i$; elements of $M^i$ are called homogeneous
	elements of degree $i$;
	\item an $I$-graded algebra is a graded module equipped with a homomorphism
	$\mu:A\otimes A\to A$, also written as a multiplication
	$xy:=\mu(x\otimes y)$, such that $A$ becomes a ring and
	\[ 1_A\in A^0,\qquad A^i\cdot A^j\subset A^{i+j},
	\qquad i,j\in I; \]
	thus the multiplication $\mu$ is completely determined by the data
	$\mu^{i,j}:A^i\otimes A^j\to A^{i+j}$;
	\item homomorphisms and isomorphisms between $I$-graded modules or algebras
	are defined in the usual way; $I$-graded subalgebras and ideals are defined
	similarly.
\end{itemize}

% segment-id: o014.li2.chapter5.multiplicative-structures.p003
For $k\in I$, also define a homomorphism of degree $k$ between $I$-graded
modules: this is a module homomorphism $\varphi:M\to N$ satisfying
$\varphi(M^i)\subset M^{i+k}$ for every $i$, and it is determined by the data
$(\varphi^i:M^i\to M^{i+k})_{i\in I}$. Taking $k=0$ recovers homomorphisms in
the usual sense. Degrees of homomorphisms add under composition.

% segment-id: o014.li2.chapter5.multiplicative-structures.d001
\begin{definition}\label{def:dg-algebra-preview}
	\index{differential graded algebra}
	Let $I$ be a commutative monoid equipped with a homomorphism
	$\epsilon:I\to\Z/2\Z$. A \emph{differential $I$-graded algebra} whose
	differential has degree $k$ is an $I$-graded algebra
	$A=\bigoplus_{p\in I}A^p$ together with an endomorphism of degree $k\in I$,
	$d=(d^p)_p:A\to A$, satisfying $d^2=0$ and the Leibniz rule
	\[ d(xy)=(dx)\cdot y+(-1)^{\epsilon(p)}x\cdot dy,
	\qquad x\in A^p,\quad y\in A^{p'},\quad p,p'\in I. \]
\end{definition}

% segment-id: o014.li2.chapter5.multiplicative-structures.p004
Observe that $d(1_A)=d(1_A\cdot1_A)=d(1_A)+d(1_A)$, hence $d(1_A)=0$.
Clearly $\Ker(d)$ is an $I$-graded subalgebra of $A$, while $\Image(d)$ is a
two-sided $I$-graded ideal in $\Ker(d)$. Therefore
\[ \Hm(A,d):=\Ker(d)/\Image(d) \]
is again an $I$-graded algebra.

% segment-id: o014.li2.chapter5.multiplicative-structures.e001
\begin{example}
	Take $I=\Z$ and $\epsilon(p)=p\;\bmod 2$. For any complex $X$ over a
	$\Bbbk$-linear category $\mathcal{B}$, the $\Hom$ complex
	$\Hom^\bullet(X,X)$ of Definition~\ref{def:Hom-cplx}, together with
	$d=d_{\Hom^\bullet(X,X)}$, becomes a differential $\Z$-graded algebra
	whose differential has degree $1$. These assertions are essentially the
	content of Lemma~\ref{prop:Hom-cplx-Leibniz}. We have
	\[ \Hm(\Hom^\bullet(X,X),d)=
	\bigoplus_{n\in\Z}\Hom_{\cate{K}(\mathcal{B})}(X,X[n]), \]
	with multiplication induced by composition of morphisms in
	$\cate{K}(\mathcal{B})$.
\end{example}

% segment-id: o014.li2.chapter5.multiplicative-structures.e002
\begin{example}
	Take $\Bbbk=\CC$, $I=\Z$, and $\epsilon(p)=p\;\bmod2$. The
	$\CC$-valued differential forms of degree $p$ on a smooth manifold
	$\mathfrak{X}$ form a vector space $A^p(\mathfrak{X})$. Set
	$A(\mathfrak{X}):=\bigoplus_{p=0}^{\dim X}A^p(\mathfrak{X})$. With
	multiplication $\mu(\omega\otimes\eta):=\omega\wedge\eta$ and exterior
	differentiation $\dd$, this becomes a differential $\Z$-graded algebra
	whose differential has degree $1$. The associated algebra
	$\Hm(A(\mathfrak{X}),\dd)=\bigoplus_p\Hm^p_{\mathrm{dR}}(\mathfrak{X})$
	is precisely the de Rham cohomology of $\mathfrak{X}$; the multiplication
	induced by $\wedge$ gives it a graded-algebra structure. This is a
	fundamental object in topology and is isomorphic as a graded algebra to the
	singular cohomology ring of $\mathfrak{X}$.

	The multiplication on $\Hm(A(\mathfrak{X}),\dd)$ also satisfies
	$xy=(-1)^{pq}yx$, where $x$ and $y$ are homogeneous elements of degrees
	$p$ and $q$, respectively, because the multiplication on
	$A(\mathfrak{X})$ already has this property. As a simple and beautiful
	example, complex projective space $\mathfrak{X}:=\mathbb{P}^n(\CC)$ gives
	the graded algebra
	$\Hm(A(\mathfrak{X}),\dd)\simeq\CC[t]/(t^{n+1})$, where the variable $t$
	corresponds to a homogeneous element of degree $2$.
\end{example}

% segment-id: o014.li2.chapter5.multiplicative-structures.p005
Returning to our main subject, in this section we consider the following
cases:
\begin{center}\begin{tabular}{|c|c|c|} \hline
	Abbreviation & $I$ & $\epsilon:I\to\Z/2\Z$ \\ \hline
	Singly graded & $\Z$ & $\epsilon(p)=p\;\bmod2$ \\
	Bigraded & $\Z^2$ & $\epsilon(p,q)=p+q\;\bmod2$ \\ \hline
\end{tabular}\end{center}

% segment-id: o014.li2.chapter5.multiplicative-structures.p006
This theory overlaps with Definition~\ref{def:differential-object}, but the
focus of this section is the multiplication, which was not mentioned there.

% segment-id: o014.li2.chapter5.multiplicative-structures.p007
To lighten the notation, below we collectively call these various cases
``graded'', distinguishing them mainly by the superscript $p$ or $(p,q)$.
Likewise, we refer collectively to differential $I$-graded algebras as
differential graded algebras.

% segment-id: o014.li2.chapter5.multiplicative-structures.c001
\begin{convention}
	For a cohomological bigraded spectral sequence $\mathscr{E}$ in the case
	$\mathcal{A}=\Bbbk\dcate{Mod}$, we identify every page $E_r$ with the
	graded module $\bigoplus_{p,q}E_r^{p,q}$ and regard
	$d_r=(d_r^{p,q})_{p,q}$ as an endomorphism of degree $(r,-r+1)$ of the
	graded module $E_r$. The homological and singly graded cases are treated
	similarly.
\end{convention}

% segment-id: o014.li2.chapter5.multiplicative-structures.p008
In short, a spectral sequence with a multiplicative structure is a spectral
sequence consisting of differential graded algebras.

% segment-id: o014.li2.chapter5.multiplicative-structures.d002
\begin{definition}\label{def:mult-ss}
	\index{spectral sequence!multiplicative structure}
	A \emph{multiplicative structure} on a cohomological bigraded spectral
	sequence $\mathscr{E}$ is a family of homomorphisms
	$\mu_r:E_r\otimes E_r\to E_r$ (that is, ``multiplications'') such that
	\begin{itemize}
		\item every $(E_r,\mu_r,d_r)$ is a differential graded algebra whose
		differential has degree $(r,-r+1)$;
		\item the map $t_{r+1}:\Hm(E_r,d_r)\rightiso E_{r+1}$ in the spectral
		sequence data is an isomorphism of graded algebras.
	\end{itemize}
	Multiplicative structures on singly graded or homological spectral sequences
	are defined similarly.
\end{definition}

% segment-id: o014.li2.chapter5.multiplicative-structures.p009
The preceding discussion shows that $Z_r=\Ker(d_r)$ is a graded subalgebra of
$E_r$, while $B_r=\Image(d_r)$ is a two-sided graded ideal of $Z_r$. Hence
$\Hm(E_r,d_r)$ becomes a graded algebra; this is the precise meaning of
Definition~\ref{def:mult-ss}.

% segment-id: o014.li2.chapter5.multiplicative-structures.p010
Furthermore, if the limit exists, then
$Z_\infty=\bigoplus_{p,q}Z_\infty^{p,q}$ is also a graded subalgebra, with
$B_\infty=\bigoplus_{p,q}B_\infty^{p,q}$ as its graded ideal. Thus
$E_\infty$ is also a graded algebra.

% segment-id: o014.li2.chapter5.multiplicative-structures.p011
\index{dg-algebra}
As an example, consider differential graded algebras whose differential has
degree $1$, abbreviated as dg-algebras. Expanding the definition, this is
equivalent to a complex of $\Bbbk$-modules
$X=(X^n,d_X^n)_{n\in\Z}$ such that
$X\xlongequal{\text{identified}}\bigoplus_nX^n$ is equipped with a
multiplication $\mu:X\otimes X\to X$ satisfying
$X^p\cdot X^{p'}\subset X^{p+p'}$ and the Leibniz rule
\[ d(xy)=dx\cdot y+(-1)^px\cdot dy,
	\qquad x\in X^p,\quad y\in X^{p'}. \]
Now suppose the data $(X,d)$ are extended to a filtered complex
$(X,d,\mathrm{F}^\bullet X)$. Recall that this already entails
$d(\mathrm{F}^pX)\subset\mathrm{F}^pX$; we further require the multiplication
on $X$ to be compatible with the filtration:
\[ \mathrm{F}^pX^n\cdot\mathrm{F}^{p'}X^{n'}
	\subset\mathrm{F}^{p+p'}X^{n+n'}. \]
In this case $(X,d,\mu,\mathrm{F}^\bullet X)$ is called a filtered
differential graded algebra. For the induced filtration on $\Hm(X,d)$, the
condition above ensures that
\[ \gr\Hm(X,d)\xlongequal{\text{identified}}
	\bigoplus_{(p,q)\in\Z^2}\gr^p\Hm^{p+q}(X,d) \]
naturally becomes a graded algebra.
\index{differential graded algebra!filtered}

% segment-id: o014.li2.chapter5.multiplicative-structures.t001
\begin{proposition}\label{prop:dga-mult-ss}
	Let $(X,d,\mu,\mathrm{F}^\bullet X)$ be a filtered differential graded
	algebra whose differential has degree $1$.
	\begin{enumerate}[(i)]
		\item The spectral sequence $\mathscr{E}$ of the filtered complex
		$(X,d,\mathrm{F}^\bullet X)$ has a canonical multiplicative structure.
		\item The isomorphisms
		$E_\infty^{p,q}\simeq\gr^p\Hm^{p+q}(X,d)$ in the classical convergence
		theorem~\ref{prop:classical-convergence} in fact give an isomorphism of
		graded algebras $E_\infty\simeq\gr\Hm(X,d)$.
	\end{enumerate}
\end{proposition}
\begin{proof}
	Take $x\in E_r^{p,q}$ and $y\in E_r^{p',q'}$. By the description in
	Proposition~\ref{prop:filtered-cplx-ss}, choose lifts
	\begin{gather*}
		\widetilde{x}\in\mathrm{F}^pX^{p+q}\cap
		d^{-1}\left(\mathrm{F}^{p+r}X^{p+q+1}\right),\qquad
		\widetilde{y}\in\mathrm{F}^{p'}X^{p'+q'}\cap
		d^{-1}\left(\mathrm{F}^{p'+r}X^{p'+q'+1}\right).
	\end{gather*}
	By definition,
	\begin{gather*}
		\widetilde{x}\widetilde{y}\in
		\mathrm{F}^{p+p'}X^{p+q+p'+q'},\\
		d(\widetilde{x}\widetilde{y})
		=d\widetilde{x}\cdot\widetilde{y}
		+(-1)^{p+q}\widetilde{x}\cdot d\widetilde{y}
		\in\mathrm{F}^{p+p'+r}X^{p+q+p'+q'+1}.
	\end{gather*}
	Thus $\widetilde{x}\widetilde{y}$ determines an element
	$xy\in E_r^{p+p',q+q'}$. A routine calculation (please verify it) shows
	that $xy$ depends only on $x$ and $y$.

	Next, $d_r:E_r\xrightarrow{(r,-r+1)}E_r$ is induced by the differential
	$d$ on $X$. This gives $E_r$ the structure of a differential graded
	algebra; associativity, the Leibniz rule, and all other required properties
	reduce to checks on $(X,d)$. The same argument shows that
	$\Hm(E_r,d_r)\simeq E_{r+1}$ is an isomorphism of graded algebras.

	The assertion that the canonical isomorphisms
	$E_\infty^{p,q}\simeq\gr^p\Hm^{p+q}(X,d)$ in the classical convergence
	theorem preserve multiplication likewise reduces to a check on $X$, so we
	omit the details.
\end{proof}

% segment-id: o014.li2.chapter5.multiplicative-structures.r001
\begin{remark}
	\index{Cartan--Eilenberg system}
	Multiplicative structures are an indispensable tool in topology. Algebraic
	topology without multiplication is almost unimaginable, or at least
	insipid. Historically, when Leray first encountered spectral sequences, he
	already considered multiplicative structures, calling them ``spectral
	rings''. They often greatly simplify spectral-sequence calculations in
	topology. This is why the present section gives a brief introduction to
	multiplicative structures, though it only scratches the surface.

	There is an evident difficulty in applying Definition~\ref{def:mult-ss}.
	The differential graded algebra structures on all the $E_r$ must be given
	at once, since the definition alone provides no reason for the
	multiplication on one page to induce a multiplication on the next. This can
	indeed be done in some situations, as in
	Proposition~\ref{prop:dga-mult-ss}.

	For general cases, or for those that are difficult to compute by hand,
	spectral sequences are often constructed from simple data, such as the exact
	couples of \S\ref{sec:exact-couples}. It is therefore natural to ask whether
	there are simple conditions on an exact couple ensuring that the corresponding
	spectral sequence has a multiplicative structure.

	The answer appears to be negative: information beyond the exact couple is
	required. One useful construction is the equally classical
	\emph{Cartan--Eilenberg system}, together with its spectral product; see
	\cite[II.A]{Dou58}. As an application, the Serre--Atiyah--Hirzebruch
	spectral sequence, which is of central importance in topology, has a
	multiplicative structure for every multiplicative generalized cohomology
	theory. Since this material departs from our main line, we stop here.
\end{remark}
